\section*{Capítulo \ref{ch:g12:meiosis-genetic-shuffling} --- La meiosis y la mezcla genética}

\begin{solution}{exo:g12:meiosis-genetic-shuffling:1}
\omterm{def:g12:meiosis-genetic-shuffling:meiosis}{Diploide}: dos copias de cada \omterm{def:g11:cell-cycle-mitosis:chromosome}{cromosoma} ($2n$); \omterm{def:g12:meiosis-genetic-shuffling:meiosis}{haploide}: una ($n$). La
\omterm{def:g12:meiosis-genetic-shuffling:meiosis}{meiosis} reduce a la mitad el número de \omterm{def:g10:universal-dna:information}{cromosomas}, de $2n$ a $n$, y
baja el \omterm{def:g10:universal-dna:information}{ADN} de $2q$ (tras la replicación) a $q/2$ por gameto.
\end{solution}

\begin{solution}{exo:g12:meiosis-genetic-shuffling:2}
La \omterm{def:g12:meiosis-genetic-shuffling:meiosis}{meiosis} I separa los \omterm{def:g10:universal-dna:information}{cromosomas} homólogos de cada pareja; la \omterm{def:g12:meiosis-genetic-shuffling:meiosis}{meiosis}
II separa las dos \omterm{def:g11:cell-cycle-mitosis:chromosome}{cromátidas} de cada \omterm{def:g11:cell-cycle-mitosis:chromosome}{cromosoma}. La primera es la
reduccional.
\end{solution}

\begin{solution}{exo:g12:meiosis-genetic-shuffling:3}
Tras la \omterm{def:g12:meiosis-genetic-shuffling:meiosis}{meiosis} I: 4 \omterm{def:g10:universal-dna:information}{cromosomas} (de dos \omterm{def:g11:cell-cycle-mitosis:chromosome}{cromátidas} cada uno), \omterm{def:g10:universal-dna:information}{ADN} $q$.
Tras la \omterm{def:g12:meiosis-genetic-shuffling:meiosis}{meiosis} II: 4 \omterm{def:g10:universal-dna:information}{cromosomas} (de una \omterm{def:g11:cell-cycle-mitosis:chromosome}{cromátida} cada uno), \omterm{def:g10:universal-dna:information}{ADN}
$q/2$.
\end{solution}

\begin{solution}{exo:g12:meiosis-genetic-shuffling:4}
$2^4 = 16$.
\end{solution}

\begin{solution}{exo:g12:meiosis-genetic-shuffling:5}
Un cruzamiento con una pareja que solo lleva \omterm{def:g10:universal-dna:gene}{alelos} \omterm{def:g11:genes-and-disease:genetic}{recesivos}, de modo
que cada descendiente muestra el gameto del progenitor puesto a prueba.
La proporción de recombinantes revela si dos \omterm{def:g10:universal-dna:gene}{genes} están en el mismo
\omterm{def:g11:cell-cycle-mitosis:chromosome}{cromosoma} y, si lo están, a qué distancia.
\end{solution}

\begin{solution}{exo:g12:meiosis-genetic-shuffling:6}
\omterm{def:g12:meiosis-genetic-shuffling:meiosis}{Meiosis} I: $2q$; entre las divisiones: $q$; gameto: $q/2$. Sin una fase
S intermedia, la segunda división vuelve a reducir el \omterm{def:g10:universal-dna:information}{ADN} a la mitad, y
eso es lo que deja al gameto con la mitad del contenido de una \omterm{def:g10:cells-common-unit:cell}{célula}
\omterm{def:g12:meiosis-genetic-shuffling:meiosis}{diploide}.
\end{solution}

\begin{solution}{exo:g12:meiosis-genetic-shuffling:7}
Cuatro clases iguales: no están ligados, están en \omterm{def:g10:universal-dna:information}{cromosomas} distintos.
Gametos $AB$, $Ab$, $aB$, $ab$, una cuarta parte cada uno.
\end{solution}

\begin{solution}{exo:g12:meiosis-genetic-shuffling:8}
Ligados: dos clases grandes ($AB$, $ab$, parentales) y dos pequeñas
($Ab$, $aB$, recombinantes). Recombinantes: $85/1000 = 8.5\%$.
\end{solution}

\begin{solution}{exo:g12:meiosis-genetic-shuffling:9}
Solo $AB$ y $ab$: un intercambio por encima de los dos \omterm{def:g10:universal-dna:gene}{genes} cambia
segmentos que no llevan ninguno de ellos, así que las combinaciones de
$A/a$ y $B/b$ no cambian. La recombinación entre dos \omterm{def:g10:universal-dna:gene}{genes} exige un
intercambio entre ellos.
\end{solution}

\begin{solution}{exo:g12:meiosis-genetic-shuffling:10}
Alrededor de 0,8, 2,9 y 16 por 1000: un factor de 20 entre los 25 y
los 42 años.
\end{solution}

\begin{solution}{exo:g12:meiosis-genetic-shuffling:11}
En la \omterm{def:g12:meiosis-genetic-shuffling:meiosis}{meiosis} I, los dos homólogos van a una misma \omterm{def:g10:cells-common-unit:cell}{célula}: dos gametos
con dos copias y dos sin ninguna. En la \omterm{def:g12:meiosis-genetic-shuffling:meiosis}{meiosis} II, las dos \omterm{def:g11:cell-cycle-mitosis:chromosome}{cromátidas}
de un \omterm{def:g11:cell-cycle-mitosis:chromosome}{cromosoma} van a una misma \omterm{def:g10:cells-common-unit:cell}{célula}: un gameto con dos copias, otro
sin ninguna y dos normales.
\end{solution}

\begin{solution}{exo:g12:meiosis-genetic-shuffling:12}
La recombinación es proporcional a la distancia: un 2\% significa que
los \omterm{def:g10:universal-dna:gene}{genes} están muy cerca, así que rara vez cae un intercambio entre
ellos; un 48\% significa que están muy lejos, con un intercambio entre
ellos en casi todas las \omterm{def:g12:meiosis-genetic-shuffling:meiosis}{meiosis}. Como un solo intercambio da un 50\% de
gametos recombinantes, los \omterm{def:g10:universal-dna:gene}{genes} ligados y lejanos se recombinan con
tanta libertad como los independientes.
\end{solution}

\begin{solution}{exo:g12:meiosis-genetic-shuffling:13}
$AB$, $Ab$, $aB$, $ab$, una cuarta parte cada uno, por segregación
independiente de las dos parejas. El \omterm{prop:g12:meiosis-genetic-shuffling:crossingover}{entrecruzamiento} no cambia nada
para estos \omterm{def:g10:universal-dna:gene}{genes}: recombina \omterm{def:g10:universal-dna:gene}{alelos} a lo largo de un \omterm{def:g11:cell-cycle-mitosis:chromosome}{cromosoma}, y estos
están en \omterm{def:g10:universal-dna:information}{cromosomas} distintos.
\end{solution}

\begin{solution}{exo:g12:meiosis-genetic-shuffling:14}
Cada progenitor pasa a cada hijo uno de dos \omterm{def:g10:universal-dna:gene}{alelos} en cada \omterm{def:g10:universal-dna:gene}{gen}, elegido
al azar; en un \omterm{def:g10:universal-dna:gene}{gen} dado, dos hermanos recibieron el mismo \omterm{def:g10:universal-dna:gene}{alelo} de un
progenitor con probabilidad $1/2$. A lo largo de muchos \omterm{def:g10:universal-dna:gene}{genes}, la
fracción compartida promedia $1/2$, pero para una pareja concreta de
hermanos varía en torno a ese valor.
\end{solution}

\begin{solution}{exo:g12:meiosis-genetic-shuffling:15}
Clones de la planta madre, genéticamente idénticos a ella y entre sí.
La planta gana una multiplicación rápida y segura de un \omterm{def:g11:enzymes-and-phenotype:phenotype}{genotipo} que
funciona aquí y ahora; pierde la variedad que permitiría a algunos
descendientes sobrevivir a un cambio de ambiente o a un parásito nuevo.
\end{solution}

\begin{solution}{pb:g12:meiosis-genetic-shuffling:1}
\textbf{1.} Hembra $b^+b$, $vg^+vg$, con $b^+vg^+$ en un \omterm{def:g11:cell-cycle-mitosis:chromosome}{cromosoma} y
$b\,vg$ en el otro (de sus dos padres puros). Macho $bb$, $vg\,vg$.

\textbf{2.} Solo gametos $b\,vg$: el macho aporta únicamente \omterm{def:g10:universal-dna:gene}{alelos}
\omterm{def:g11:genes-and-disease:genetic}{recesivos}, así que el \omterm{def:g11:enzymes-and-phenotype:phenotype}{fenotipo} de cada descendiente revela el gameto de
la hembra.

\textbf{3.} $b^+vg^+$ (gris-normal), $b\,vg$ (negro-vestigial),
$b^+vg$ (gris-vestigial), $b\,vg^+$ (negro-normal).

\textbf{4.} Gris-normal 42\%, negro-vestigial 41\%, gris-vestigial
9\%, negro-normal 8\%.

\textbf{5.} En el mismo \omterm{def:g11:cell-cycle-mitosis:chromosome}{cromosoma}: dos clases grandes y dos pequeñas en
lugar de $1:1:1:1$.

\textbf{6.} Gris-normal y negro-vestigial: las combinaciones que
llevaban sus dos homólogos, heredadas intactas de sus padres puros.

\textbf{7.} Gris-vestigial y negro-normal, de un \omterm{prop:g12:meiosis-genetic-shuffling:crossingover}{entrecruzamiento} entre
los dos \omterm{def:g10:universal-dna:gene}{genes} en la profase I.

\textbf{8.} $(206 + 185)/2300 \approx 17\%$. Un \omterm{prop:g12:meiosis-genetic-shuffling:crossingover}{entrecruzamiento} entre
los \omterm{def:g10:universal-dna:gene}{genes} deja recombinantes dos de las cuatro \omterm{def:g11:cell-cycle-mitosis:chromosome}{cromátidas}, así que un
17\% de gametos recombinantes significa un intercambio en alrededor del
34\% de las \omterm{def:g12:meiosis-genetic-shuffling:meiosis}{meiosis}.

\textbf{9.} Un solo intercambio produce a la vez las dos \omterm{def:g11:cell-cycle-mitosis:chromosome}{cromátidas}
recombinantes, una de cada clase; y los dos homólogos van a los gametos
por igual, así que las dos clases parentales también coinciden.

\textbf{10.} Las clases parentales serían ahora gris-vestigial y
negro-normal, con un 41,5\% cada una; y las recombinantes, gris-normal
y negro-vestigial, con un 8,5\% cada una.

\textbf{11.} Sí: recombinación muy por debajo del 50\% con los dos.

\textbf{12.} $b$ --- $cn$ --- $vg$: 9 y 8 suman 17.

\textbf{13.} La probabilidad de un intercambio en un intervalo es
proporcional a su longitud, así que las frecuencias de intervalos
contiguos se suman: la frecuencia de recombinación mide la distancia a
lo largo del \omterm{def:g11:cell-cycle-mitosis:chromosome}{cromosoma}.

\textbf{14.} Un intercambio deja dos \omterm{def:g11:cell-cycle-mitosis:chromosome}{cromátidas} recombinantes de cuatro:
como mucho un 50\%; los intercambios adicionales vuelven a barajar,
pero nunca suben la fracción recombinante por encima de la mitad.

\textbf{15.} Está en otro \omterm{def:g11:cell-cycle-mitosis:chromosome}{cromosoma}, o muy lejos en el mismo; los
cruzamientos no pueden decir cuál de las dos cosas.

\textbf{16.} $2^4 = 16$ gametos; $16 \times 16 = 256$ combinaciones.

\textbf{17.} Cada intercambio produce \omterm{def:g10:universal-dna:information}{cromosomas} que no existían antes,
en una posición que varía de una \omterm{def:g12:meiosis-genetic-shuffling:meiosis}{meiosis} a otra; el número de
\omterm{def:g10:universal-dna:information}{cromosomas} distintos y, por tanto, de gametos, no tiene límite
práctico.

\textbf{18.} $2^{23} \times 2^{23} \approx \num{7e13}$: setecientas
veces más combinaciones que el número de humanos que han vivido, y eso
antes del \omterm{prop:g12:meiosis-genetic-shuffling:crossingover}{entrecruzamiento}.

\textbf{19.} Los hermanos sacan sus \omterm{def:g10:universal-dna:gene}{alelos} de los mismos cuatro juegos
(dos por progenitor) y comparten, de media, la mitad de ellos; los
desconocidos los sacan de juegos distintos.

\textbf{20.} El 17\%, la fracción de gametos de la hembra recombinados
entre $b$ y $vg$, mide la distancia entre los \omterm{def:g10:universal-dna:gene}{genes}; la segregación
independiente de las parejas y el \omterm{prop:g12:meiosis-genetic-shuffling:crossingover}{entrecruzamiento} dentro de cada
pareja hacen distinto cada gameto.
\end{solution}
